% Paper 2 (SNN-WebGPU) Supplemental Material
% Moved from main paper 2026-05-23 for NeurIPS 9-page main-text budget.
% Not counted in the 9-page limit; submit as separate supplemental PDF.
\documentclass[10pt]{article}
\usepackage[margin=1in]{geometry}
\usepackage{times}
\usepackage[utf8]{inputenc}
\usepackage[T1]{fontenc}
\usepackage{amsmath,amssymb}
\usepackage{booktabs}
\usepackage{enumitem}
\usepackage{hyperref}
\newcommand{\snnwebgpu}{\texttt{@holoscript/snn-webgpu}}
\newcommand{\cael}{CAEL}
\newcommand{\measuredFrom}[1]{\footnote{Measured from \texttt{#1}.}}

\begin{document}
\title{SNN-WebGPU: Supplemental Material}
\maketitle

\appendix

\section{Twin Test Equivalence}
\label{app:twin-test}

To satisfy the digital-twin reproducibility envelope (W.315), we maintain a bit-exact CPU reference implementation of the LIF neuron model (\texttt{src/poc/cpu-reference.ts}) and assert cross-platform parity against the WebGPU compute shader on every commit.
The reference uses the identical mathematical model---exponential decay, threshold-and-reset, and a $2\,$ms refractory period---implemented in plain TypeScript with deterministic xorshift32 stimulus generation so the same seed produces the same synaptic-current pattern on every backend.

\begin{table}[h]
\centering
\caption{Twin-test parity: CPU reference vs. WebGPU LIF kernel.
  Measured 2026-05-04 on NVIDIA Ampere (RTX 3060 Laptop, driver 591.74) via \texttt{vitest} with the Dawn-native WebGPU backend.
  Tolerance: $5\times10^{-5}$ absolute / $10^{-4}$ relative on membrane potentials; spike masks exact ($0$ vs. $1$).}
\label{tab:twin-test}
\small
\measuredFrom{HoloScript/packages/snn-webgpu/src/paper/__tests__/LIFTwinTest.test.ts}%
\begin{tabular}{@{}rrrr@{}}
\toprule
\textbf{Neurons} & \textbf{Ticks} & \textbf{Max abs. diff.} & \textbf{Spike parity} \\
\midrule
1{,}024   & 100 & $\!<\!5\times10^{-5}$ & 1{,}024 / 1{,}024 \\
65{,}536  & 10  & $\!<\!5\times10^{-5}$ & 65{,}536 / 65{,}536 \\
2{,}048   & 20  & $\!<\!5\times10^{-5}$ & 2{,}048 / 2{,}048 (custom $\tau{=}10$, $v_{\text{thresh}}{=}{-}50$) \\
\bottomrule
\end{tabular}
\end{table}

Table~\ref{tab:twin-test} reports the assertion outcomes for three canonical configurations.
The residual floating-point divergence ($\sim$10$^{-5}$) is bounded by IEEE-754 f32 \texttt{exp()} differences between the WGSL shader and the JS host; spike masks remain exact because the threshold comparison is a discrete predicate that is insensitive to sub-ulp variation.
This twin-test assertion is the empirical substrate for the cross-backend determinism claim in the evaluation section of the main paper: if CPU and GPU agree at the f32 level, observed cross-GPU hash mismatches (if any) point to driver or compiler bugs, not model drift. The twin-test is the program's W.315 wire-format equivalent primitive specialized to spike data: the WGSL shader (one implementation) and the JS host (a second implementation) constitute a two-implementation agreement asserted by deterministic replay through the SNN shadow-replay harness at \texttt{packages/core/src/paper-0c-spike/stage-c-shadow-replay.ts}, which compares the runtime spike digest against the canonical JSONL spike digest per step (the \texttt{digest\_match} bit-identity check) and falls through to bounded-loss tolerance ($q_f/2$ for $\geq 99.99\%$ of steps) when the cross-backend regime would relax bit-identity at sub-ulp scale.

\section{Decoder Cost Analysis}
\label{app:decoder-cost}

For every ``novel / faster / better'' claim, we publish the decoder (readback) cost on the same row (W.314).
The LIF kernel writes membrane potentials and spike flags to GPU storage buffers; consuming them requires a GPU-to-CPU readback.
Table~\ref{tab:decoder-cost} quantifies this overhead across the neuron counts used in the throughput evaluation.

\begin{table}[h]
\centering
\caption{GPU-to-CPU readback cost for membrane-potential and spike-mask buffers.
  Measured 2026-05-04 on NVIDIA Ampere (RTX 3060 Laptop, driver 591.74) via \texttt{vitest} with the Dawn-native WebGPU backend.
  Simulation cost is the compute-only time for the same 100-tick run; readback is serial (worst-case) after the final tick.
  Effective throughput is bytes transferred divided by readback latency.}
\label{tab:decoder-cost}
\small
\measuredFrom{HoloScript/packages/snn-webgpu/src/paper/__tests__/LIFDecoderCost.test.ts}%
\begin{tabular}{@{}rrrrrr@{}}
\toprule
\textbf{Neurons} & \textbf{Readback (ms)} & \textbf{Sim. (ms)} & \textbf{Ratio} & \textbf{$\mu$s/tick} & \textbf{Throughput} \\
\midrule
1{,}024   & 0.73  & 22.7  & 3.1\%  & 7.3  & 0.01 GB/s \\
4{,}096   & 0.59  & 20.3  & 2.8\%  & 5.9  & 0.06 GB/s \\
16{,}384  & 0.67  & 18.1  & 3.6\%  & 6.7  & 0.19 GB/s \\
65{,}536  & 0.89  & 23.0  & 3.7\%  & 8.9  & 0.59 GB/s \\
262{,}144 & 2.09  & 27.3  & 7.1\%  & 20.9 & 1.00 GB/s \\
\bottomrule
\end{tabular}
\end{table}

The readback cost is $O(N)$ in neuron count and remains below $8\%$ of the compute cost across the evaluated range.
At the canonical operating point ($N{=}65{,}536$, 100 ticks), the incremental per-tick readback is $8.9\;\mu$s---negligible compared to the $230\;\mu$s per-tick simulation time.
The asymptotic bottleneck is therefore the LIF kernel itself, not the decoder.
Because the WebGPU \texttt{mapAsync} path on Chromium copies through a shared CPU/GPU memory mapping on integrated GPUs and through PCIe on discrete GPUs, the measured numbers are an upper bound: platforms with unified memory (Apple Silicon, Intel UHD) exhibit lower readback latency than the discrete-GPU figures reported here.

\section{User Study Protocol Details}
\label{app:user-study-protocol}
\label{supp:user-study}

\paragraph{Design.}
The experiment uses a within-subjects (repeated-measures) design with
two tool conditions:
\begin{itemize}[leftmargin=*,topsep=2pt]
  \item \textbf{Condition A} — the \snnwebgpu{} Studio viewer with
    embodied-navigation replay, per-tick CAEL-trace inspection, and
    a toggleable tropical-activation overlay.
  \item \textbf{Condition B} — baseline: raw spike-raster CSV, a
    PyTorch CPU reference script, and standard matplotlib visualisation.
    No provenance trace or formal-bridge overlay is available.
\end{itemize}
Order is counterbalanced with a Latin square; scenario difficulty is
fixed (easy $\to$ medium $\to$ hard) within each condition block.

\paragraph{Participants.}
Target sample: pilot $N=3$ (descriptive only), replication $N=12$
(inferential).  Recruitment is from the HoloScript community Discord,
the neuromorphic-computing sub-reddit, and a local university CS/NEU
mailing list.  Inclusion criteria: age $\geq 18$, self-rated neural-network
familiarity $\geq 3$ on a 1--5 scale, prior use of a Python ML
framework, and a WebGPU-capable browser.  Exclusions: authors or
direct collaborators on \snnwebgpu{}, participants who have previously
seen the stimulus scenarios, and individuals with colour-blindness that
impedes reading the navigation-world colour coding (red=obstacle,
green=goal).  Compensation is a \$30 USD gift card.  IRB approval is
pending; this preregistration is committed to the repository before any
external recruitment (D.011 reproducibility gate).

\paragraph{Materials and apparatus.}
The three diagnostic scenarios are drawn from the embodied navigation
task described in the main paper.  Each scenario presents a failed
$128$-neuron LIF run in a $20 \times 20$ room with $6$ obstacles:
\begin{enumerate}[leftmargin=*,topsep=2pt]
  \item \emph{Obstacle overshoot} (easy): the agent collides with an
    obstacle near $(10,10)$ and never recovers.  The correct diagnosis
    is an under-inhibited east-direction sensor population.
  \item \emph{Goal miss} (medium): the agent reaches the goal
    neighbourhood ($(16.8,17.9)$) but circles without touching the
    target at $(18,18)$.  The correct diagnosis is insufficient goal
    excitation relative to residual obstacle inhibition.
  \item \emph{Sensor-noise drift} (hard): Gaussian noise
    ($\sigma=0.3$) on a single diagonal distance sensor causes a
    compounding drift into a side wall over $80$ ticks.  The correct
    diagnosis requires identifying the noisy sensor from anomalous
    spike-rate patterns.
\end{enumerate}
In Condition A the participant clicks through the replay, inspects CAEL
trace entries, and optionally toggles the tropical overlay to see which
directional population was under-inhibited.  In Condition B the participant
scrolls the CSV and cross-references the trajectory plot against the
reference script's encoding logic.

\paragraph{Measures.}
\begin{itemize}[leftmargin=*,topsep=2pt]
  \item \textbf{Diagnosis accuracy} — binary, coded against a rubric
    by two independent researchers (Cohen's $\kappa$ reported).
  \item \textbf{Time to diagnosis} — seconds from stimulus onset to
    confidence-submit, auto-timestamped by the instrument.
  \item \textbf{Confidence} — 7-point Likert (``How confident are you
    that your diagnosis is correct?''); primary continuous DV for
    inferential testing.
  \item \textbf{Cognitive load} — single-item NASA-TLX mental-demand
    rating (1--7), chosen for brevity at small $N$.
  \item \textbf{Perceived usefulness} — 7-point Likert (``How useful
    would this tool be for your own neuromorphic research?''),
    collected once per condition block.
  \item \textbf{SUS} — 10-item System Usability Scale, applied to
    each tool condition separately.
  \item \textbf{CAEL-specific trust} (Condition A only) — 7-point
    Likert: ``How confident are you that the CAEL trace would detect
    tampering if someone altered a spike-rate vector mid-run?''
\end{itemize}

\paragraph{Procedure.}
Each participant attends a single $35$--$45$ minute session.  After
informed consent and demographics ($3$ min), the participant completes
the three scenarios in Condition A or B ($15$ min), answers the
post-condition SUS and free-response debrief ($2$ min), switches to the
second condition and repeats the three scenarios ($15$ min), and
concludes with a preference question, open-ended comparison, and
demographics recap ($3$ min).  The debrief script discloses the study's
purpose and offers a summary of findings on request.

\paragraph{Analysis plan.}
The primary hypothesis ($H_1$) states that Condition A increases
diagnosis confidence relative to Condition B.  Because the data are
ordinal Likert scores and $N \leq 12$, we use the Wilcoxon
signed-rank test (paired by participant) at $\alpha = 0.05$
(two-tailed).  Effect size is reported as rank-biserial correlation
$r$ with Cliff's delta as a robustness check.  Secondary DVs (time,
cognitive load, usefulness) are corrected with Holm-Bonferroni.  Pilot
($N=3$) results are descriptive only (means/medians with 95\% CI,
explicit ``not inferential'' disclaimer).  Exclusions: participants who
complete fewer than four of six trials; response-set participants
(flagged and discussed before dropping).  Stopping rule is fixed $N$;
no interim peeking.

\paragraph{Expected outcomes and claim rung.}
We expect Condition A to outperform Condition B on confidence and
accuracy, with the largest effect on the hard (sensor-noise) scenario
where CAEL anomaly indicators and tropical overlays provide the most
leverage over raw CSV inspection.  This submission contributes the
committed pre-data bundle only (D.011 evidence surface); the paper
upgrades to rung~2 after IRB-approved pilot execution and to rung~3
if the replication completes before camera-ready.  The preregistration,
task script, analysis template, and raw-data specification are committed
to \texttt{research/user-study-data/paper-2/}.

% ============================================================================
\section{Implementation Details (Full)}
\label{supp:implementation}

\subsection{Software Stack}

\snnwebgpu{} is implemented in TypeScript (strict ES2020) with WGSL compute
shaders, built with \texttt{tsup} and tested with \texttt{vitest}.
The package structure:

\begin{itemize}
\item \texttt{gpu-context.ts} --- WebGPU adapter/device lifecycle, capability detection,
  maximum buffer size negotiation (default 256MB).
\item \texttt{buffer-manager.ts} --- GPU buffer allocation with staging pools for
  async readback, 4-byte alignment enforcement.
\item \texttt{pipeline-factory.ts} --- Shader module compilation and pipeline caching;
  6 shader modules $\times$ 16 entry points.
\item \texttt{lif-simulator.ts} --- Single-population LIF simulator; manages
  5 GPU buffers (params, membrane, synaptic input, spikes, refractory).
\item \texttt{snn-network.ts} --- Multi-layer network orchestrator with named layers,
  configurable connections, and optional STDP per connection.
\item \texttt{spike-codec.ts} --- \texttt{SpikeEncoder} and \texttt{SpikeDecoder}
  classes mapping continuous data $\leftrightarrow$ spike trains.
\item \texttt{traits/TropicalActivationTrait.ts} --- SNN--ReLU bridge (CPU path);
  GPU path via \texttt{tropical-activation.wgsl}.
\item \texttt{graph/TropicalShortestPaths.ts} --- APSP and SSSP via tropical
  GEMM and SpMV, with automatic CPU/GPU routing based on problem size.
\end{itemize}

\subsection{The SNNCognitionEngine}

The \texttt{SNNCognitionEngine} in the engine package bridges \snnwebgpu{} to
the \cael{} agent loop.
It implements the \texttt{CAELCognitionEngine} interface:

\begin{enumerate}
\item \textbf{Sensor-to-current mapping}: Physical sensor readings (Pa, K, m/s$^2$)
  are normalized to $[0, 1]$ per batch, then scaled to $[0, I_\mathrm{scale}]$ mV
  (default: 20mV).
  This injection range drives neurons from $V_\mathrm{rest} = -65$ mV toward
  $V_\theta = -55$ mV: 20\% input yields light spiking; 100\% yields near-saturation.

\item \textbf{Multi-step LIF}: Each \texttt{think()} call runs
  $\lceil \Delta t \cdot 1000 / \Delta t_\mathrm{LIF} \rceil$ LIF steps,
  collecting all spikes with millisecond-resolution timestamps.

\item \textbf{Goal derivation}: Aggregate firing rate drives a simple GOAP goal stack:
  $>30\%$ firing $\Rightarrow$ ``stabilize'',
  $>5\%$ $\Rightarrow$ ``monitor'',
  else $\Rightarrow$ ``idle''.

\item \textbf{GPU/CPU fallback}: The engine attempts WebGPU initialization; if
  unavailable, it falls back to a bit-exact \texttt{CPUReferenceSimulator}
  implementing the same LIF equations on the CPU.
\end{enumerate}

\subsection{The LIF Kernel in Detail}

The core LIF compute shader processes one neuron per thread:

\begin{verbatim}
@compute @workgroup_size(256)
fn lif_step(@builtin(global_invocation_id) gid: vec3<u32>) {
  let i = gid.x;
  if (i >= params.neuron_count) { return; }
  var v = membrane_v[i];
  let i_syn = synaptic_input[i];
  var ref_timer = refractory[i];
  let decay = exp(-params.dt / params.tau);
  if (ref_timer > 0.0) {
    refractory[i] = max(ref_timer - params.dt, 0.0);
    membrane_v[i] = params.v_reset;
    spikes[i] = 0.0;
    return;
  }
  v = params.v_rest + (v - params.v_rest) * decay + i_syn;
  if (v >= params.v_threshold) {
    spikes[i] = 1.0;
    membrane_v[i] = params.v_reset;
    refractory[i] = 2.0; // 2ms refractory
  } else {
    spikes[i] = 0.0;
    membrane_v[i] = v;
  }
}
\end{verbatim}

A batch variant (\texttt{lif\_step\_batch}) processes 4 sub-steps per dispatch
for higher temporal resolution when latency is less critical than throughput.

\subsection{STDP Learning}

The STDP kernel implements soft-bounded Hebbian plasticity:
\begin{equation}
\Delta W_{ji} = \begin{cases}
\eta (1 - W_{ji}) & \text{if } s_i = 1 \text{ and } s_j = 1 \quad \text{(LTP)} \\
-0.5\eta W_{ji}  & \text{if } s_i = 1 \text{ and } s_j = 0 \quad \text{(LTD)} \\
0                 & \text{otherwise}
\end{cases}
\end{equation}
Weights are clamped to $[0, 1]$ after each update.
The soft bounds prevent weight explosion (LTP pushes toward 1 with diminishing force)
and ensure weight decay (LTD pulls toward 0).

\subsection{Wait-Free Parallel Sorting}

For spike train processing requiring sorted order, we implement a wait-free
Blelloch prefix-sum radix sort in WGSL (\texttt{blelloch-scan-sort.wgsl}).
The algorithm avoids spin-waits that deadlock on Apple/Intel GPUs:
Phase 1 computes local histograms per workgroup (shared memory atomics);
Phase 2 runs a Blelloch up-sweep/down-sweep scan for global prefix sums;
Phase 3 scatters elements to sorted positions.
A single-block fast path handles arrays $\leq 256$ elements entirely in
shared memory.

\subsection{Test Coverage}

The \snnwebgpu{} package is validated by Vitest suites covering
LIF simulation, spike codecs, SNN graph wiring, tropical kernels (including APSP
benchmarks), GPU buffer/pipeline paths, and integration scenarios.
A full-package \texttt{vitest run} on 2026-04-17 reported \textbf{205} passing tests
(Windows/Node; some suites exercise native WebGPU when available, others use mocks).

% ============================================================================
\section{Related Work (Full)}
\label{supp:related}

\paragraph{SNN Simulation Frameworks.}
NEST~\cite{gewaltig2007nest} is the gold standard for large-scale SNN simulation
($>10^6$ neurons) but requires a Linux cluster.
BrainCog~\cite{tang2024braincog} provides a PyTorch-based SNN framework with
embodied AI integration (iCub, Robosim) but requires CUDA GPUs.
snnTorch~\cite{eshraghian2023training} focuses on training SNNs with surrogate
gradients---a complementary concern to our GPU inference focus.
SpikingJelly~\cite{fang2023spikingjelly} offers CUDA-accelerated SNN simulation
with neuromorphic dataset support.
None of these run in a web browser.

\paragraph{Web-Based Neural Networks.}
TensorFlow.js~\cite{smilkov2019tensorflow} supports DNNs on WebGPU but does
not implement spiking dynamics.
ONNX Runtime Web uses WebGPU for standard neural network inference.
Neither provides spike encoding/decoding, STDP, or LIF neuron models.
To our knowledge, \snnwebgpu{} is the first SNN implementation targeting WebGPU.

\paragraph{Browser-Native SNN Prior Work and Comparison Scope.}
The throughput figures reported in this paper (peak $1.48 \times 10^9$ and
$5.89 \times 10^9$ neuron-timesteps/s on integrated and discrete WebGPU,
respectively) are scoped to \emph{browser-native, zero-installation deployment}
and should not be compared directly against CUDA-backed frameworks such as
SpikingJelly or BrainCog, which exploit dedicated GPU driver paths unavailable
in the browser sandbox.
The meaningful comparison class for \snnwebgpu{} is prior browser-accessible SNN work.
The closest precursor is the WebGL-based SNN-P systems simulator~\cite{cabarle2022webgl},
which runs spiking neural P systems in web browsers using the WebGL rasterization pipeline.
That approach is limited to WebGL's compute-hostile rasterization model and does
not support STDP, multi-layer LIF dynamics, or hash-chain provenance.
\snnwebgpu{} replaces the WebGL rasterization pipeline with WebGPU compute
shaders---providing true compute parallelism, read/write storage buffers, and
workgroup-local shared memory---enabling biologically realistic LIF neurons,
Hebbian plasticity, and cryptographic audit trails, all within the browser
sandbox with no CUDA driver required.

\paragraph{Neuromorphic Hardware.}
Intel Loihi 2~\cite{davies2018loihi} and IBM TrueNorth~\cite{merolla2014truenorth}
implement SNNs in custom silicon with $>10^6$ neurons and $< 1$W power.
SpiNNaker~\cite{furber2014spinnaker} uses ARM cores for flexible neuron models.
These provide superior energy efficiency but require specialized hardware
procurement and toolchains---the antithesis of our zero-installation approach.

\paragraph{Spiking World Models.}
Capone and Paolucci~\cite{capone2024dreaming} (Sci.\ Rep.\ 2024) demonstrate
biologically plausible model-based reinforcement learning in recurrent spiking
networks, where ``dreaming'' accelerates skill acquisition without offline replay
buffers.
Our work complements this by showing that spike-based world-model execution can
run end-to-end in a browser with provenance over every spike, rather than
requiring HPC infrastructure or specialized neuromorphic hardware.

\paragraph{Tropical Algebra in Neural Networks.}
The connection between tropical geometry and neural networks has been explored
theoretically~\cite{zhang2018tropical, alfarra2022decision}, showing that ReLU
networks partition input space into tropical hypersurfaces.
Maragos et al.~\cite{maragos2021tropical} provide a comprehensive survey of
tropical algebra in signal processing and machine learning.
Our contribution is the \emph{implementation}: a GPU-accelerated tropical bridge
that connects SNN rate codes to ReLU activations in a working system with
empirical benchmarks.

\paragraph{Provenance in Simulation.}
Simulation provenance has been explored in scientific workflows~\cite{simmhan2005survey}
and digital twins~\cite{tao2019digital}, but always at the simulation level---not
at the individual neural decision level.
Our \cael{} system provides per-spike-train, per-action hash-chain provenance,
enabling counterfactual analysis with cryptographic integrity guarantees.

% ============================================================================
\section{Background: Extended Derivations}
\label{supp:background}

\subsection{LIF Model: Full Equations}

The membrane potential $V(t)$ of neuron $i$ evolves as:
\begin{equation}
V_i(t + \Delta t) = V_\mathrm{rest} + \bigl(V_i(t) - V_\mathrm{rest}\bigr) \cdot e^{-\Delta t / \tau} + I_\mathrm{syn}^{(i)}(t)
\end{equation}
where $\tau$ is the membrane time constant, $V_\mathrm{rest}$ is the resting potential,
and $I_\mathrm{syn}^{(i)}(t)$ is the total synaptic input current.
When $V_i(t + \Delta t) \geq V_\theta$ (the spike threshold), the neuron emits a spike,
resets to $V_\mathrm{reset}$, and enters a refractory period $t_\mathrm{ref}$.
The synaptic current for post-synaptic neuron $j$ is:
\begin{equation}
I_\mathrm{syn}^{(j)}(t) = \sum_{i} W_{ji} \cdot s_i(t)
\end{equation}
where $W_{ji}$ is the synaptic weight from neuron $i$ to $j$, and $s_i(t) \in \{0, 1\}$
is the spike indicator.

\subsection{Tropical Semiring Axioms}

The max-plus semiring $(\mathbb{R} \cup \{-\infty\}, \max, +)$ satisfies:
$a \oplus b = \max(a, b)$; $a \otimes b = a + b$; identity for $\oplus$: $-\infty$;
identity for $\otimes$: $0$.
The min-plus semiring $(\mathbb{R} \cup \{+\infty\}, \min, +)$ satisfies:
$a \oplus b = \min(a, b)$; identity: $+\infty$ (encoded as \texttt{1e30}).
Repeated squaring $A^{\otimes \lceil \log_2 N \rceil}$ yields APSP distances.
See also Appendix~B (Tropical Semiring Properties) of the main paper.

\end{document}
